\documentclass{article}
\usepackage{graphicx}
\usepackage{xeCJK}
\usepackage{fontspec}
\usepackage{amsmath}

\title{ML part}
\author{saiisagod}
\date{2026年1月31日}

\begin{document}

\maketitle

\noindent\textbf{说在前面：}
本文档面经来自于小红书、牛客网、以及各个仓库整理所得，真假混杂谨慎使用(╹ڡ╹ )。

\section{WXG}
\begin{enumerate}
\item \textbf{看了很多帖子，个人感觉WXG不喜欢问八股喜欢专注于项目本身。}\\
\item \textbf{prenorm和postnorm，grpo损失}\\
欠拟合指模型过于简单，无法捕捉数据的真实模式，在训练集和测试集上表现都不好；过拟合指模型过于复杂，过度记忆训练集噪声，在训练集表现很好但测试集性能很差。解决欠拟合可以：增加模型复杂度（更复杂模型、增加特征）、减少正则强度、训练更久。解决过拟合可以：增加数据（数据增强）、使用正则化（L1/L2）、Dropout、简化模型结构、使用交叉验证选择合适的超参数、早停等。

\item \textbf{场景题：设置一个模型预测用户停留时间}\\
欠拟合：模型学不出规律，训练集表现就差（loss高/准确率低），测试集也差，二者都不好、差距不大。过拟合：模型把训练集“记住了”，训练集表现很好（loss低/准确率高），但测试集明显变差，训练与测试差距很大。

\item \textbf{最长递增子序列}\\
学习算法的期望泛化误差可分解为：误差 = 偏差$^2$ + 方差 + 不可约误差。偏差反映模型表达能力是否足够（偏差大对应欠拟合），方差反映模型对训练数据扰动的敏感程度（方差大对应过拟合）。提高模型复杂度通常会降低偏差、增加方差，因此需要在偏差和方差之间取得平衡，通过模型选择和正则化实现。

\item \textbf{了解不了解推荐系统，传统推荐和生成式推荐的区别，大模型能替代推荐系统的哪个阶段}\\
典型现象是：训练集误差很低，但验证集/测试集误差明显偏高，两者差距较大。使用学习曲线（训练集大小 vs 误差）也能判断：若训练误差很低但测试误差不随数据量增加明显下降，则为高方差。应对方式包括：减小模型复杂度、增加正则化、增加训练数据、使用集成方法等。

\item \textbf{1、空间O(1)时间O(n)找数组中重复数据 2、数组中第k大的数字}\\
深度模型参数量巨大，函数空间极其丰富，有能力拟合训练数据中的噪声和偶然结构；如果数据量相对不足、正则化不充分，就容易出现训练误差很小而测试误差较大的情况。此外，深网对输入扰动较敏感（如对抗样本），也体现了高方差特性。
\item \textbf{最长不重复子序列，有效括号}\\
\item \textbf{最回文子串}\\
\end{enumerate}

\end{document}
